% ============================================================
% Supplementary Material — DocVerify (PRL submission)
% Companion document archived on Zenodo. Contains the full
% experimental tables that did not fit the 7-page PRL limit.
% Build: pdflatex supplementary.tex  (x2)
% ============================================================
\documentclass[11pt]{article}

\usepackage[a4paper,margin=2.3cm]{geometry}
\usepackage[T1]{fontenc}
\usepackage[utf8]{inputenc}
\usepackage{amsmath}
\usepackage{amssymb}
\usepackage{booktabs}
\usepackage{array}
\usepackage{multirow}
\usepackage{graphicx}
\usepackage{xcolor}
\usepackage{caption}
\usepackage[colorlinks=true,linkcolor=blue,citecolor=blue,urlcolor=blue]{hyperref}

\captionsetup{font=small,labelfont=bf}
\renewcommand{\arraystretch}{1.15}

% Supplementary numbering: S1, Table S1, etc.
\renewcommand{\thesection}{S\arabic{section}}
\renewcommand{\thetable}{S\arabic{table}}
\renewcommand{\thefigure}{S\arabic{figure}}

\title{\textbf{Supplementary Material}\\[4pt]
\large Multi-objective Pareto Hyperparameter Optimization for Joint
Detection and Localization of Forged Identity Documents}
\author{Hern\'{a}n D\'{i}az Rodr\'{i}guez\\
\small Department of Computer Science, University of Oviedo, Asturias, Spain\\
\small \texttt{diazhernan@uniovi.es}}
\date{}

\begin{document}
\maketitle

\noindent
This document is a companion to the Pattern Recognition Letters submission
\emph{Multi-objective Pareto Hyperparameter Optimization for Joint Detection
and Localization of Forged Identity Documents} (DocVerify). It collects the
complete experimental tables that, owing to the journal's strict page limit,
are summarized in only one or two sentences in the main paper or in the
response to reviewers. All numbers are computed from the released result CSVs
(archived together with this document); section references (e.g.\ Sec.~4.3)
point to the main paper. Unless stated otherwise, all blind-test results use
the held-out 15\% FantasyID partition, $224\times224$ inputs, and the
Pareto-selected configuration of Table~2 in the main paper. Statistical tests
are paired two-sided Wilcoxon signed-rank tests with Holm correction across
metrics within an experiment; effect sizes are paired Cohen's~$d$.

\tableofcontents
\vspace{6pt}
\hrule

% ============================================================
\section{Computational cost breakdown}
\label{sec:compute}

The ${\sim}59$-hour budget reported in Sec.~4.1 is the cost of the full
\emph{scientific} protocol (unbiased nested cross-validation plus a 30-seed
ablation), not the cost of \emph{using} the method. Table~\ref{tab:compute}
gives the measured wall-clock breakdown on a single NVIDIA RTX~5060\,Ti
(16\,GB) with VRAM caching. Once the configuration is fixed, retraining the
single deployed model from scratch (up to 100 epochs with early stopping)
costs only $7.7\pm1.0$~min per seed; the Pareto loss weighting adds no
inference-time cost.

\begin{table}[h]
  \centering
  \caption{Wall-clock cost per phase (single GPU, VRAM-cached). Per-seed
           retraining time is the mean over the 30 blind-test seeds
           (range 6.0--10.2\,min).}
  \label{tab:compute}
  \begin{tabular}{lrr}
    \toprule
    Phase & Trainings & Wall-clock \\
    \midrule
    Nested HPO search (10 outer $\times$ 50 MOTPE $\times$ 5 inner) & 2\,500 & ${\sim}50$\,h \\
    Outer final models (one retrain per outer fold)                & 10     & ${\sim}1$\,h  \\
    Blind-test ablation (30 seeds $\times$ 4 variants)             & 120    & ${\sim}12$\,h \\
    \midrule
    \textbf{Total scientific protocol}                             & \textbf{2\,630} & \textbf{${\sim}59$--$63$\,h} \\
    \midrule
    \textbf{Practical deployment} (one fixed config, one seed)     & \textbf{1} & \textbf{$7.7\pm1.0$\,min} \\
    \bottomrule
  \end{tabular}
\end{table}

% ============================================================
\section{Data-augmentation ablation}
\label{sec:aug}

Sec.~4.1 states that standard augmentation degraded every metric on the
FantasyID holdout. Table~\ref{tab:aug} gives the full 30-seed comparison.
The augmented model applies, in the training loader only, random rotation
($\pm10^\circ$), brightness, contrast, Gaussian blur, and JPEG compression,
using the \emph{same} Pareto configuration. On the clean synthetic holdout,
augmentation significantly degrades every metric (all Holm $p<10^{-7}$) and
inflates variance by $2$--$5\times$, confirming the paper's no-augmentation
choice for this dataset.

\begin{table}[h]
  \centering
  \caption{Augmentation ablation (blind test, 30 seeds). $\Delta$:
           aug $-$ no\_aug; $d$: paired Cohen's~$d$; $p_{\text{Holm}}$:
           Holm-corrected paired Wilcoxon. All differences significant.}
  \label{tab:aug}
  \begin{tabular}{lccccc}
    \toprule
    Metric & No aug (ours) & Aug & $\Delta$ & $d$ & $p_{\text{Holm}}$ \\
    \midrule
    PR-AUC    & $0.9967\pm0.0021$ & $0.9577\pm0.0457$ & $-0.0390$ & $-0.84$ & $1.5\times10^{-8}$ \\
    Dice      & $0.8750\pm0.0180$ & $0.7797\pm0.0396$ & $-0.0953$ & $-2.10$ & $9.3\times10^{-9}$ \\
    BAcc      & $0.9573\pm0.0184$ & $0.8302\pm0.1007$ & $-0.1271$ & $-1.20$ & $5.6\times10^{-9}$ \\
    F1$_1$    & $0.9651\pm0.0152$ & $0.8593\pm0.1061$ & $-0.1058$ & $-0.97$ & $1.1\times10^{-8}$ \\
    F1$_{\text{macro}}$ & $0.9423\pm0.0234$ & $0.7998\pm0.1136$ & $-0.1425$ & $-1.19$ & $7.5\times10^{-9}$ \\
    \bottomrule
  \end{tabular}
\end{table}

% ============================================================
\section{Class-imbalance (\texttt{pos\_weight}) sensitivity}
\label{sec:posweight}

Sec.~3.1 notes that the mild class imbalance (2.5:1, attack-majority) leaves
PR-AUC robust and that a \texttt{pos\_weight} sweep confirms no significant
effect. Table~\ref{tab:posweight} reports the three settings (inverse-frequency
$0.397$, neutral $1.0$, and the attack-favouring $2.518$), each over 30 seeds.
No metric differs significantly from the neutral baseline after Holm correction
(all $p_{\text{Holm}}>0.5$), so no loss re-weighting is applied in the main
model.

\begin{table}[h]
  \centering
  \caption{\texttt{pos\_weight} sensitivity (blind test, 30 seeds).
           Baseline is \texttt{pos\_weight}$=1.0$. No comparison against the
           baseline reaches significance after Holm correction.}
  \label{tab:posweight}
  \begin{tabular}{lccccc}
    \toprule
    \texttt{pos\_weight} & PR-AUC & Dice & BAcc & Rec.\ attack & Rec.\ bonafide \\
    \midrule
    $0.397$ (inv.\ freq.)        & $0.9973\pm0.0018$ & $0.8793\pm0.0168$ & $0.9614\pm0.0146$ & $0.9431$ & $0.9797$ \\
    $1.0$ (baseline)             & $0.9967\pm0.0021$ & $0.8750\pm0.0180$ & $0.9573\pm0.0184$ & $0.9435$ & $0.9712$ \\
    $2.518$ (attack-favouring)   & $0.9967\pm0.0019$ & $0.8688\pm0.0240$ & $0.9552\pm0.0164$ & $0.9435$ & $0.9669$ \\
    \bottomrule
  \end{tabular}
\end{table}

% ============================================================
\section{Selection criterion: Pareto vs.\ scalar (full results)}
\label{sec:scalar}

Table~3 of the main paper reports a compact form of the selection-criterion
study. Table~\ref{tab:scalar} gives the complete 30-seed results with 95\%
confidence intervals for all five metrics, and Table~\ref{tab:scalar_stats}
the paired statistics. Both scalar criteria select the global argmax (by
validation PR-AUC or validation Dice) over the \emph{same} 500 MOTPE trials;
our Pareto criterion uses the distance to the ideal point. Pareto improves
localization over \texttt{by\_prauc} and improves \emph{both} objectives over
\texttt{by\_dice}.

\begin{table}[h]
  \centering
  \caption{Selection-criterion comparison (blind test, 30 seeds).
           Mean\,$\pm$\,Std with 95\% CI in brackets.}
  \label{tab:scalar}
  \small
  \begin{tabular}{lccc}
    \toprule
    Metric & \textbf{Pareto (ours)} & \texttt{by\_prauc} & \texttt{by\_dice} \\
    \midrule
    PR-AUC & $0.9967\pm0.0021$        & $0.9953\pm0.0021$        & $0.9948\pm0.0027$ \\
           & {\scriptsize[0.9959, 0.9975]} & {\scriptsize[0.9945, 0.9961]} & {\scriptsize[0.9938, 0.9958]} \\
    Dice   & $0.8750\pm0.0180$        & $0.8589\pm0.0230$        & $0.8566\pm0.0147$ \\
           & {\scriptsize[0.8683, 0.8818]} & {\scriptsize[0.8504, 0.8675]} & {\scriptsize[0.8511, 0.8620]} \\
    BAcc   & $0.9573\pm0.0184$        & $0.9481\pm0.0146$        & $0.9463\pm0.0170$ \\
           & {\scriptsize[0.9504, 0.9642]} & {\scriptsize[0.9427, 0.9536]} & {\scriptsize[0.9400, 0.9527]} \\
    F1$_1$ & $0.9651\pm0.0152$        & $0.9569\pm0.0122$        & $0.9564\pm0.0121$ \\
           & {\scriptsize[0.9594, 0.9708]} & {\scriptsize[0.9523, 0.9614]} & {\scriptsize[0.9519, 0.9609]} \\
    F1$_{\text{macro}}$ & $0.9423\pm0.0234$ & $0.9294\pm0.0184$   & $0.9283\pm0.0184$ \\
           & {\scriptsize[0.9335, 0.9510]} & {\scriptsize[0.9225, 0.9362]} & {\scriptsize[0.9214, 0.9352]} \\
    \bottomrule
  \end{tabular}
\end{table}

\begin{table}[h]
  \centering
  \caption{Paired statistics, Pareto vs.\ each scalar criterion (30 seeds).
           $\Delta$: Pareto $-$ scalar; $d$: paired Cohen's~$d$;
           $p_{\text{Holm}}$: Holm-corrected paired Wilcoxon;
           \checkmark{}: significant at $p_{\text{Holm}}<0.05$.}
  \label{tab:scalar_stats}
  \small
  \begin{tabular}{llcccc}
    \toprule
    Comparison & Metric & $\Delta$ & $d$ & $p_{\text{Holm}}$ & sig. \\
    \midrule
    \multirow{5}{*}{Pareto vs \texttt{by\_prauc}}
      & PR-AUC & $+0.0014$ & $0.49$ & $0.062$ & \\
      & Dice   & $+0.0161$ & $0.59$ & $0.043$ & \checkmark \\
      & BAcc   & $+0.0092$ & $0.40$ & $0.061$ & \\
      & F1$_1$ & $+0.0082$ & $0.45$ & $0.066$ & \\
      & F1$_{\text{macro}}$ & $+0.0129$ & $0.46$ & $0.050$ & \checkmark \\
    \midrule
    \multirow{5}{*}{Pareto vs \texttt{by\_dice}}
      & PR-AUC & $+0.0019$ & $0.64$ & $0.0072$ & \checkmark \\
      & Dice   & $+0.0185$ & $1.00$ & $0.00039$ & \checkmark \\
      & BAcc   & $+0.0110$ & $0.45$ & $0.040$ & \checkmark \\
      & F1$_1$ & $+0.0087$ & $0.44$ & $0.066$ & \\
      & F1$_{\text{macro}}$ & $+0.0140$ & $0.47$ & $0.043$ & \checkmark \\
    \bottomrule
  \end{tabular}
\end{table}

% ============================================================
\section{Architecture-agnostic validation (ResNet-18)}
\label{sec:resnet}

Sec.~4.6 states that the identical MOTPE protocol, applied to an
ImageNet-pretrained ResNet-18 backbone, yields comparable blind-test
performance and that neither model dominates. Table~\ref{tab:resnet} gives the
full 30-seed comparison. The ResNet-18 Pareto-selected configuration
(distance to ideal $=0.164$) is $\eta=3.68\times10^{-4}$,
$\lambda_{\text{wd}}=1.45\times10^{-6}$, dropout $=0.381$,
$C_{\text{dec}}=256$, $\lambda_{\text{mask}}=2.69$, the last consistent with
DocVerify's $2.46$. DocVerify is better on PR-AUC; ResNet-18 is better on Dice
and the F1 metrics; balanced accuracy does not differ significantly.

\begin{table}[h]
  \centering
  \caption{DocVerify vs.\ ResNet-18 under the identical MOTPE protocol
           (blind test, 30 seeds). $\Delta$: DocVerify $-$ ResNet-18;
           $d$: paired Cohen's~$d$; $p_{\text{Holm}}$: Holm-corrected paired
           Wilcoxon.}
  \label{tab:resnet}
  \begin{tabular}{lccccc}
    \toprule
    Metric & DocVerify & ResNet18-MOTPE & $\Delta$ & $d$ & $p_{\text{Holm}}$ \\
    \midrule
    PR-AUC    & $0.9967\pm0.0021$ & $0.9942\pm0.0025$ & $+0.0024$ & $+0.75$ & $2.8\times10^{-4}$ \\
    Dice      & $0.8750\pm0.0180$ & $0.8894\pm0.0111$ & $-0.0144$ & $-0.60$ & $8.1\times10^{-3}$ \\
    BAcc      & $0.9573\pm0.0184$ & $0.9629\pm0.0108$ & $-0.0055$ & $-0.21$ & $0.285$ (n.s.) \\
    F1$_1$    & $0.9651\pm0.0152$ & $0.9779\pm0.0069$ & $-0.0128$ & $-0.67$ & $3.8\times10^{-3}$ \\
    F1$_{\text{macro}}$ & $0.9423\pm0.0234$ & $0.9613\pm0.0114$ & $-0.0190$ & $-0.63$ & $4.0\times10^{-3}$ \\
    \bottomrule
  \end{tabular}
\end{table}

% ============================================================
\section{Cross-domain generalization (SIDTD)}
\label{sec:sidtd}

Sec.~4.6 reports that the FantasyID-trained model is near chance zero-shot on
the out-of-domain SIDTD dataset, but recovers strong performance when the same
architecture is retrained on the SIDTD training partition. Table~\ref{tab:sidtd}
gives the full curve from zero-shot ($n_{\text{shots}}=0$) through few-shot to
full retraining ($n_{\text{shots}}=1777$). Few-shot adaptation up to 400
samples stays near chance; only full in-domain retraining recovers performance,
confirming that the gap reflects domain shift rather than architectural
capacity.

\begin{table}[h]
  \centering
  \caption{SIDTD adaptation curve. $n_{\text{shots}}=0$ is zero-shot
           (30 seeds); $n_{\text{shots}}\geq10$ are few-shot/full retrainings
           (25 seeds each). Templates subset (1\,000 bonafide, 1\,222 attack).}
  \label{tab:sidtd}
  \begin{tabular}{rcccc}
    \toprule
    $n_{\text{shots}}$ & PR-AUC & ROC-AUC & BAcc & F1$_{\text{macro}}$ \\
    \midrule
    0 (zero-shot) & $0.5554\pm0.0110$ & $0.5041\pm0.0094$ & $0.5010$ & --- \\
    10   & $0.5485\pm0.0084$ & $0.4893\pm0.0147$ & $0.5002\pm0.0135$ & $0.4841\pm0.0199$ \\
    25   & $0.5595\pm0.0089$ & $0.4996\pm0.0092$ & $0.5046\pm0.0146$ & $0.4755\pm0.0255$ \\
    50   & $0.5575\pm0.0102$ & $0.4960\pm0.0128$ & $0.4953\pm0.0167$ & $0.4652\pm0.0207$ \\
    100  & $0.5706\pm0.0208$ & $0.5060\pm0.0199$ & $0.5032\pm0.0203$ & $0.4848\pm0.0259$ \\
    200  & $0.5774\pm0.0170$ & $0.5089\pm0.0144$ & $0.5031\pm0.0182$ & $0.4853\pm0.0289$ \\
    400  & $0.5912\pm0.0238$ & $0.5123\pm0.0171$ & $0.5028\pm0.0184$ & $0.4843\pm0.0349$ \\
    \textbf{1\,777 (full)} & $\mathbf{0.8780\pm0.0161}$ & $\mathbf{0.8263\pm0.0228}$ & $\mathbf{0.7188\pm0.0349}$ & $\mathbf{0.7142\pm0.0333}$ \\
    \bottomrule
  \end{tabular}
\end{table}

% ============================================================
\section{Augmentation does not close the cross-domain gap}
\label{sec:aug_sidtd}

As a follow-up to Sections~\ref{sec:aug} and~\ref{sec:sidtd}, we tested whether
augmentation-trained checkpoints generalize better zero-shot to SIDTD.
Table~\ref{tab:aug_sidtd} shows a null result: both the augmented and
non-augmented models remain at chance, with no significant difference in
detection quality after Holm correction. Augmentation therefore does not
substitute for in-domain data; closing the gap requires multi-source training
(future work, Sec.~6).

\begin{table}[h]
  \centering
  \caption{Zero-shot SIDTD generalization, augmented vs.\ non-augmented
           FantasyID training (30 seeds each). $\Delta$: aug $-$ no\_aug;
           $p_{\text{Holm}}$: Holm-corrected paired Wilcoxon.}
  \label{tab:aug_sidtd}
  \begin{tabular}{lcccc}
    \toprule
    Metric & No aug & Aug & $\Delta$ & $p_{\text{Holm}}$ \\
    \midrule
    PR-AUC  & $0.5554\pm0.0110$ & $0.5639\pm0.0096$ & $+0.0085$ & $0.077$ (n.s.) \\
    ROC-AUC & $0.5041\pm0.0094$ & $0.5089\pm0.0060$ & $+0.0048$ & $0.199$ (n.s.) \\
    \bottomrule
  \end{tabular}
\end{table}

\vspace{8pt}
\noindent\textbf{Data availability.} The result CSVs underlying every table
above, together with the experiment scripts, are released in the
reproducibility archive accompanying this document.

\end{document}
